\section{Outline} 
\label{sec:introduction:outline}

The remainder of this dissertation is organized as follows:
\begin{itemize}
    \item Chapter \ref{chap:background} (Background and Preliminaries) establishes the foundations required for this work. It introduces the principles of Bayesian inference, the mechanics of deep generative models, particularly normalizing flows, and the connection from Bayesian inference to control theory. 
    \item Chapter \ref{chap:literature} (Literature Review) provides a comprehensive survey of the current state of the art in applying deep generative models in parameter identification, state estimation, and optimal control or reinforcement learning. Through this review, the specific research gaps addressed by this dissertation are identified.
    \item Chapter \ref{chap:methods} (Methodology and Numerical Validation) details the core technical frameworks developed in this work. These methodologies are categorized into three primary research pillars: parameter identification, state estimation, and optimal control. Each method is accompanied by numerical experiments and synthetic benchmarks to provide theoretical proof of concept and demonstrate the efficacy of the proposed algorithms.
    \item Chapter \ref{chap:experiments} (Experiments and Results) presents a series of benchmarks and real-world case studies designed to evaluate the proposed frameworks. This chapter validates the performance of the developed methods in terms of accuracy, computational efficiency, and robustness to noise, following the same three-pillar structure as the methodology.
    \item Chapter \ref{chap:conclusion} (Conclusion and Future Work) synthesizes the key findings of this research, reflects on the contributions made to the fields of system identification, state estimation, and optimal control, and outlines potential directions for future research.
\end{itemize}
